\documentclass[11pt, oneside]{article} 
\usepackage[margin=3.4cm]{geometry} 
\geometry{letterpaper}
\usepackage{graphicx}
\usepackage{hyperref}
\usepackage{amssymb}
\usepackage{caption}
\usepackage{float}
\usepackage{setspace}
\usepackage{esvect}
\usepackage{url}
\usepackage{lipsum} 

%\setcounter{secnumdepth}{0}
\onehalfspacing

\newcounter{refno}
\newcommand{\reflabel}[1]{\refstepcounter{refno}\label{#1}[\arabic{refno}]}

\begin{document}
% Title Page
\begin{titlepage}
	\newcommand{\HRule}{\rule{\linewidth}{0.5mm}}
	
	\center
	\textsc{\LARGE University of Calgary}\\[1.5cm]
	\textsc{\Large ENGO 651: Advanced Geospatial Topics}\\[0.5cm]
	\textsc{\large Final Project}\\[0.5cm]

	\HRule\\[0.4cm]
	{\huge\bfseries Fuellytics: Real-Time Vehicular Fuel Consumption and Emissions Monitoring}\\[0.4cm]
	\HRule\\[1.5cm]
	
	{\large\textit{Group \#}}\\
	Adam \textsc{Smith} (30031453)\\
	Chavisa \textsc{Sornsakul} (12345678)\\
	Wai Ka \textsc{Wong Situ} (12345678)\\
	
	\vfill\vfill\vfill
	{\large\today}
	
	\vfill\vfill\vfill
	\includegraphics[width=8cm,]{img/schulich.png}\\[1cm]	
\end{titlepage}

\section{Executive summary}

\section{Problem statement}
According to the International Energy Agency, transportation accounts for almost one-quarter of global greenhouse gas emissions, and within that, road transport is responsible for the largest share of emissions \reflabel{transport-emissions}.  The burning of fossil fuels in vehicles produces carbon dioxide, which is the most prevalent greenhouse gas contributing to global warming.  Although electric vehicles are becoming more popular in the developed world, still 91\% of all transport relies on oil-based products such as gasoline, which is only a 3\% drop from the early 1970's \reflabel{oil-based-transport}.  Reducing carbon emissions from vehicles is crucial to mitigate the harmful effects of climate change and reduce the environmental burden of vehicular transport.

Tracking vehicle fuel consumption is essential to understanding the amount of carbon emissions produced by vehicles. Fuel consumption data can provide valuable insights into the efficiency of a vehicle and its impact on the environment, and can also provide insights to the driver about their driving style and fuel costs.  However, fuel consumption is heavily dependent on various factors such as vehicle type, terrain, and driving style, making it difficult to measure over short time periods.  According to the US Department of Energy, obeying speed limits, accelerating and braking gently, and reading the road ahead can improve fuel economy by up to 30\% on highways and up to 40\% in stop-and-go traffic \reflabel{driving-style}.  Moreover, engine size, vehicle weight, and driving in hilly or mountainous areas can drastically increase fuel consumption and, consequently, carbon emissions.  Some modern vehicles are equipped with a dashboard gauge that displays fuel consumption while driving, however, these gauges are often vague and uninformative, usually showing nothing more than an unlabelled bar or dial which increases while accelerating.  Without a means to directly monitor fuel consumption, drivers are often left unaware of their fuel use and how their driving style and driving conditions can reduce their vehicle's impact on the environment.

The presented smartphone application, \textit{Fuellytics}, provides an innovative solution to monitoring fuel consumption in real-time, and provides insights and reports to drivers about their fuel use while driving as well as on a weekly and monthly basis.  Fuellytics helps drivers better understand the environmental impacts of their vehicle, make educated decisions regarding transportation, and save money on fuel.

\section{Similar solutions and available literature}

Industrial tool targeted at fleet vehicles: FuelForce \url{https://info.gartnerdigitalmarkets.com/multiforce-gdm-lp?channel=capterra}, Apollo 2: \url{https://www.thetriscangroup.com/products/triscan-systems/fuel-management/}, many more...

Open tools to track fuel use, but no real time data: Fuelly: \url{https://www.fuelly.com/}, FuelLine: \url{https://fuel-line.fly.dev/}, etc.

Papers on the subject: Main reference: \url{https://www.mdpi.com/2076-3417/9/7/1369#B21-applsci-09-01369}

\section{Solution summary}

\section{Architecture}

\subsection{Design rationale}

\subsection{Architecture description}

\subsection{API}

\subsection{Sequence Diagram}

\subsection{Data models and JSON encodings}

\section{Results}

\section{Lessons learned}

\section{Conclusion}
 
 
\section{References}

\begin{small}

\noindent [\ref{transport-emissions}] \url{https://www.iea.org/reports/transport-energy-and-co2}

\noindent [\ref{oil-based-transport}] \url{https://www.iea.org/reports/transport}

\noindent [\ref{driving-style}] \url{https://afdc.energy.gov/conserve/behavior_techniques.html}


\end{small}

\end{document}